\chapter{Prompt Specifications}
\label{app:prompts}

This appendix reproduces the production prompt used by \texttt{src/prompt\_builder.py}.
Every LLM in the evaluation block in Chapter~\ref{ch:experimental-setup}
receives this prompt with a single formal bound substituted into the user
section.

\section{System Prompt}
\label{sec:system-prompt}

The system prompt is constant across every call and defines the analyst role
together with the strict rules that the explanation must respect.

\begin{lstlisting}[language=,caption={Production system prompt.}]
You are an industrial control systems network security analyst.

Your task is to explain a formally bounded IEC-104 network observation in
precise operational English.

Strict rules:
1. Use only facts explicitly provided in the formal bound.
2. Do not invent attack types, malicious intent, causality, timing patterns,
   affected assets, or mitigation actions.
3. Do not claim that the observation is confirmed normal or confirmed
   malicious unless the formal bound explicitly states that.
4. Treat proxy formal classes as observation labels, not as verified attack
   labels.
5. Stay consistent with the source identifier, destination identifier, ports,
   packet size, ASDU fields, protocol hint, and communication direction.
6. If evidence is insufficient for a security conclusion, explicitly say that
   the bound does not confirm a security conclusion.
7. Write 3 to 5 sentences.
8. Use professional analyst language.
\end{lstlisting}

\section{User Prompt Template}
\label{sec:user-prompt}

The user prompt is generated from a single formal bound. The placeholder
fields are filled in by \texttt{prompt\_builder.build\_user\_prompt}.

\begin{lstlisting}[language=,caption={User prompt template.}]
Explain the following formally bounded ICS network event.

Event metadata:
- Event ID: <event_id>
- Dataset: <dataset>
- Event granularity: <event_granularity>
- Timestamp ms: <timestamp_ms>
- Formal bound version: <formal_bound_version>
- Schema status: <schema_status>

Formal observation:
- Protocol hint: <protocol_hint>
- Communication direction: <communication_direction>
- Frame type: <frame_type>
- Observation pattern: <observation_pattern>
- Formal class: <formal_class>
- Formal class status: <formal_class_status>

Observed features:
- source_id: <source_id>
- destination_id: <destination_id>
- bytes: <bytes>
- pkt_length: <pkt_length>
- srcport: <srcport>
- dstport: <dstport>
- asdu_address: <asdu_address>
- asdu_cot: <asdu_cot>
- asdu_items: <asdu_items>
- asdu_type: <asdu_type>
- frame_fmt: <frame_fmt>

Allowed claims:
<list of ten atomic statements that may be made>

Required claims:
<list of five statements that must be made>

Forbidden claims:
<list of eight statements that must not be made>

Write the explanation now.
\end{lstlisting}

\section{Allowed, Required, and Forbidden Claims}
\label{sec:claim-lists}

The three claim lists are generated per bound by
\texttt{formal\_bound\_builder.py}. Allowed claims are atomic statements
derivable from the observed features. Required claims state that the
explanation must mention the source identifier, the destination identifier,
the port numbers, the ASDU fields, and the proxy nature of the formal class.
Forbidden claims state that the explanation must not assert a confirmed
attack, confirmed malice, an arbitrary attack class, additional devices, a
causal explanation, an unobserved timing pattern, a mitigation requirement,
or verified normal traffic.

\section{Local Model Configuration}
\label{sec:local-model-config}

Every local model is called with identical generation parameters:

\begin{itemize}
  \item Temperature: 0.1
  \item Top-p: 0.9
  \item Backend: Ollama 0.23.3 with 4-bit quantisation
  \item HTTP timeout: 180 seconds per request
\end{itemize}

The values are persisted in
\texttt{src/pipeline\_config.py} (\texttt{OLLAMA\_TEMPERATURE},
\texttt{OLLAMA\_TOP\_P}, \texttt{OLLAMA\_TIMEOUT\_SECONDS}) so that any run
can be reproduced from the configuration alone.

\section{Cloud Model Configuration}
\label{sec:cloud-model-config}

The cloud runner uses the same generation parameters where the vendor API
supports them and a maximum-token cap of 600 tokens. The defaults can be
overridden per run through environment variables
(\texttt{OPENAI\_MODEL}, \texttt{GEMINI\_MODEL},
\texttt{REASONGUARD\_CLOUD\_PILOT\_LIMIT}).

\section{Prompt Variants}
\label{sec:prompt-variants}

The exposé specifies three prompt variants to be compared in a 20-sample
pilot at the start of the evaluation block. The variants will be inserted
here once the W11 pilot is run; the production variant in
Sections~\ref{sec:system-prompt}--\ref{sec:user-prompt} is the current
default and is also the variant used to produce every result in
Chapter~\ref{ch:results}.
